% Risk-Aware Multi-Criteria Selective Prediction
% IEEE conference paper skeleton, per project_plan.md §10.
%
% NOTE: this uses the generic `IEEEtran.cls` (conference mode), fetched from
% CTAN. The project plan's submission checklist (§10) explicitly warns that
% venue-specific IEEE templates differ -- swap in the exact template for
% your target venue before submission. This environment has no LaTeX
% toolchain installed, so this file has been written carefully by hand but
% NOT compiled; compile it yourself (Overleaf or a local TeX Live install)
% before trusting it.
\documentclass[conference]{IEEEtran}

\usepackage{cite}
\usepackage{amsmath,amssymb}
\usepackage{graphicx}
\usepackage{booktabs}
\usepackage{multirow}
\usepackage{algorithm}
\usepackage{algorithmic}
\usepackage{xcolor}

\newcommand{\todo}[1]{\textcolor{red}{\textbf{[TODO: #1]}}}

\begin{document}

\title{Risk-Aware Multi-Criteria Selective Prediction}

\author{
  \IEEEauthorblockN{Yash Kale}
  \IEEEauthorblockA{\todo{Department, Institution, City, Country}\\
  Email: yash.kale@uptiq.ai}
}

\maketitle

\begin{abstract}
Selective classifiers abstain on inputs they are unlikely to classify
correctly, trading coverage for reliability. The dominant heuristic --
thresholding the base model's maximum softmax probability (Chow's rule) --
is Bayes-optimal only when the model's posteriors are correct, an
assumption broken by miscalibration, distribution shift, and epistemic
uncertainty. We study whether aggregating a heterogeneous bank of
uncertainty signals (calibrated confidence, ensemble disagreement, an
energy/OOD score, and feature-space geometry) with a learned,
coverage-targeted aggregator improves the accuracy-coverage trade-off over
the best single signal, and whether any such advantage is regime-dependent
(tied in-distribution, larger under shift and at low coverage). We further
wrap the resulting score in a distribution-free conformal risk-control
layer and audit the approach for subgroup disparity. \todo{Fill in once
Table~\ref{tab:aurc_summary} and the shift results (§9 week 9) are final --
state the headline number here, not before.}
\end{abstract}

\begin{IEEEkeywords}
selective classification, uncertainty quantification, conformal prediction, abstention, OOD detection
\end{IEEEkeywords}

\section{Introduction}
\label{sec:intro}

Chow's rule~\cite{chow1970} -- abstain when the maximum predicted
posterior falls below a threshold -- is the Bayes-optimal selective
classifier \emph{if} the model's posteriors are the true class posteriors.
They rarely are: neural and tree-based classifiers are commonly
miscalibrated~\cite{guo2017calibration}, behave unpredictably on
out-of-distribution inputs~\cite{liu2020energy}, and carry epistemic
uncertainty that a single forward pass cannot express. A large benchmark
of selective-classification methods across dozens of datasets found that
no single method dominates across operating points and
objectives~\cite{benchmark2024}, which is the central motivation for an
\emph{adaptive}, multi-signal approach rather than one more single-signal
heuristic.

At the same time, a well-tuned softmax-confidence baseline is a
notoriously strong competitor: several recent studies show that elaborate
learned-abstention methods fail to beat it
in-distribution~\cite{feng2023better}, and that cheap post-hoc fixes (
temperature scaling, logit normalization) recover much of the claimed gain
of heavier methods~\cite{cattelan2023logitnorm}. We treat this as a
falsifiable prediction rather than a threat to route around: we
pre-register (\S\ref{sec:rq}) that any advantage from aggregation is
\emph{regime-dependent} -- small or absent in-distribution, and larger
under distribution shift and at low coverage, where a single signal's
blind spots matter most.

\textbf{Contributions.}
\begin{enumerate}
  \item A post-hoc, model-agnostic multi-criteria abstention framework
    requiring no retraining of the base model.
  \item A coverage-targeted training objective for the aggregator, shown
    to \todo{outperform / tie -- report the actual A1-vs-A2 result once
    final} correctness-classification training in the low-coverage
    regime.
  \item An instance-adaptive signal weighting mechanism, with analysis of
    which signals dominate in which failure regime.
  \item Distribution-free selective-risk control via a conformal wrapper,
    evaluated across tabular and image benchmarks including distribution
    shift and subgroup disparity.
\end{enumerate}

\subsection{Research questions}
\label{sec:rq}
\textbf{RQ1.} Does aggregating heterogeneous uncertainty signals improve
the accuracy-coverage trade-off over the best single-signal abstention
rule?
\textbf{RQ2.} Is that advantage regime-dependent -- roughly tied with
confidence thresholding in-distribution, clearly better under shift and at
low coverage?
\textbf{RQ3.} Can the learned score be wrapped in a finite-sample
guarantee on selective risk without losing coverage relative to a
heuristic threshold?
\textbf{RQ4.} Does multi-criteria abstention reduce the concentration of
abstentions (and residual error) on minority subgroups documented by prior
work~\cite{jones2021disparities}?

\section{Related work}
\label{sec:related}
\todo{Expand into full prose once the gap table
(\texttt{paper/related\_work\_gap\_table.md}) is populated with verified
citations. Every entry in that table must be checked against IEEE
Xplore/arXiv/Semantic Scholar before being cited here -- see the caution
in project\_plan.md \S2.}

\subsection{Selective classification foundations}
\cite{chow1970, geifman2017selective, elyaniv2010foundations}

\subsection{Learned abstention}
\cite{geifman2019selectivenet, liu2019deepgamblers, corbiere2019confidnet, huang2020sat, mozannar2020defer, moon2020crl}

\subsection{The strong-baseline literature}
\cite{feng2023better, cattelan2023logitnorm}

\subsection{Fairness in selective classification}
\cite{jones2021disparities}

\subsection{Uncertainty and OOD signal sources}
\cite{guo2017calibration, lakshminarayanan2017ensembles, gal2016dropout, liu2020energy, lee2018mahalanobis, sun2022knn, jiang2018trustscore, huang2021gradnorm}

\subsection{Conformal prediction and risk control}
\cite{vovk2005conformal, angelopoulos2022riskcontrol, angelopoulos2021learnthentest, angelopoulos2021gentle}

\todo{Close this section with the gap table (columns: method, signals
used, needs retraining?, guarantee?, evaluated under shift?, subgroup
analysis?) as a compact figure/table -- the empty cells are the paper's
justification.}

\section{Preliminaries}
\label{sec:prelim}
Let $f$ be a fixed, frozen $K$-class base classifier and $g: \mathcal{X}
\to \{0,1\}$ a selection function; $(f,g)$ abstains on $x$ when $g(x)=0$.
\textbf{Coverage} is $\Pr[g(x)=1]$. \textbf{Selective risk} is the error
rate of $f$ restricted to the accepted region,
$R(f,g) = \mathbb{E}[\ell(f(x),y) \cdot g(x)] / \mathbb{E}[g(x)]$.
The \textbf{risk-coverage curve} plots selective risk as a function of the
target coverage achieved by thresholding a score $s(x)$; \textbf{AURC} is
its area, and \textbf{E-AURC} subtracts the AURC of the best achievable
(oracle) ordering, which controls for the confounding effect of base
accuracy when aggregating across datasets~\cite{geifman2017selective}.

\section{Method}
\label{sec:method}

\subsection{Signal bank}
Given $x$, we compute $u(x) \in \mathbb{R}^m$ from four tiers: (A)
logit-derived signals available from a single forward pass -- max softmax
probability, entropy, margin, negative energy, logit-norm-normalized MSP,
temperature-scaled MSP; (B) model-internal signals requiring multiple
forward passes or an ensemble -- vote entropy, mean pairwise KL, top-class
variance; (C) representation/data-geometry signals requiring the training
set at inference -- $k$NN distance, local label agreement, trust
score~\cite{jiang2018trustscore}, Mahalanobis distance to the predicted
class centroid~\cite{lee2018mahalanobis}. \todo{Add Tier D
(tree-specific) signals if/when GBM-internal disagreement signals are
implemented -- currently out of scope, see PROJECT\_STATUS.md.}

\subsection{Aggregator}
We compare three aggregator classes of increasing expressiveness: A0, a
learning-free rank/z-score average; A1, logistic regression or LightGBM
stacked on $u(x)$ to predict $\mathbf{1}[f(x)\neq y]$ (ordinary stacking);
and A2, the same function class as A1 trained instead with a
coverage-targeted loss (\S\ref{sec:loss}).

\subsection{Coverage-targeted loss}
\label{sec:loss}
Let $g(x) = \sigma((\tau - s(x))/T)$ be a soft acceptance gate. We
minimize
\begin{equation}
\mathcal{L}_\kappa = \frac{\sum_i g(x_i)\,\ell(f(x_i),y_i)}{\sum_i g(x_i)}
  + \lambda \max(0, \kappa - \tfrac{1}{n}\textstyle\sum_i g(x_i))^2,
\end{equation}
with $\tau$ a learned threshold, which trains the aggregator to rank well
\emph{at the coverage the system is actually deployed at} rather than for
correctness-classification accuracy overall. We also report an AURC
surrogate (summing $\mathcal{L}_\kappa$ over a coverage grid) and a
pairwise ranking loss (logistic loss on $s(\text{incorrect}) -
s(\text{correct})$ margins).

\subsection{Instance-adaptive signal weighting}
A gating network $w(x) = \mathrm{softmax}(h(u(x)))$ produces
per-instance weights over the (standardized) signals, $s(x) =
\sum_j w_j(x)\, \tilde u_j(x)$, so the model can route to different
signals for different failure regimes (e.g. Mahalanobis distance for
far-OOD inputs, local label agreement for ambiguous in-distribution ones).

\subsection{Conformal risk-control wrapper}
On a held-out calibration set of size $n$, we select the largest threshold
$\hat\tau$ such that a finite-sample upper confidence bound on selective
risk stays below a user target $\alpha$, giving
$\Pr[R(f, g_{\hat\tau}) \le \alpha] \ge 1-\delta$ under exchangeability,
independent of how the aggregator was trained. This guarantee is valid
in-distribution; under distribution shift, exchangeability fails and the
guarantee is evaluated empirically rather than claimed.

\section{Experimental setup}
\label{sec:setup}
\todo{Datasets (Table~\ref{tab:datasets}), base models, splits and
cross-fitting, baselines, metrics, statistical protocol -- fill in from
\texttt{results/results.parquet} once the full 12-dataset, 10-seed sweep
(project\_plan.md \S9 week 8) is complete. Current results cover
\{adult, german\_credit\} at \{1, 5\} seeds -- see PROJECT\_STATUS.md for
exact scope.}

\section{Results}
\label{sec:results}
\todo{Main AURC/E-AURC table, risk-coverage figures, coverage-at-
guaranteed-risk table, shift results, cost curves, CD diagram. Every
number here must be generated by \texttt{scripts/make\_tables.py} /
\texttt{scripts/make\_figures.py} querying \texttt{results/results.parquet}
-- never hand-typed.}

\section{Ablations and analysis}
\label{sec:ablations}
\todo{Leave-one-signal-out; signal-family-only; aggregator class; loss
comparison (A1-vs-A2, the headline ablation for the "this is just
stacking" objection); with/without conformal wrapper; naive-vs-cross-fit
meta-fitting; meta-set-size sweep; signal correlation/PCA.}

\section{Limitations}
\label{sec:limitations}
The conformal guarantee (\S\ref{sec:method}) assumes exchangeability and
is therefore valid in-distribution only; Tier C signals require the
training set at inference time, adding memory and latency cost; the
meta-split needed to train the aggregator is itself a cost on small
datasets (quantified by the meta-set-size ablation); and current results
are tabular-and-CIFAR-scale, not ImageNet-scale.

\section{Conclusion}
\label{sec:conclusion}
\todo{Write last, after Results is final.}

\section*{Acknowledgment}
\todo{Remove or anonymize for double-blind submission per the checklist in
project\_plan.md \S10.}

\bibliographystyle{IEEEtran}
\bibliography{references}

\end{document}
